% SOURCE_STATEMENT: "For entanglement entropy in disordered free-fermion systems, how much delocalization of the one-particle eigenvalues and eigenfunctions is needed to produce a logarithmic enhancement of the area law, even when the spectrum is not absolutely continuous?"
% Preparation and final self-audit: CONTRIBUTING.md prompts C and D; source/audit artifacts retained in the local physics_statements_2026_10_11 batch.
% !TeX program = lualatex
\documentclass[11pt]{article}
\usepackage[margin=25mm]{geometry}
\usepackage{fontspec}
\setmainfont{Latin Modern Roman}
\usepackage{amsmath,amssymb,mathtools}
\usepackage{unicode-math}
\setmathfont{Latin Modern Math}
\usepackage{xurl,hyperref}
\hypersetup{colorlinks=true,urlcolor=blue}
\setcounter{secnumdepth}{0}
\setlength{\parindent}{0pt}
\setlength{\parskip}{5pt}
\setlength{\emergencystretch}{3em}
\newcommand{\sref}[2]{\hyperlink{source:#1:#2}{[S#2]}}
\newcommand{\cataloguescope}[1]{\paragraph{Recorded target.}#1}
\newcommand{\R}{\mathbb R}
\newcommand{\C}{\mathbb C}
\newcommand{\N}{\mathbb N}
\newcommand{\Z}{\mathbb Z}
\newcommand{\Q}{\mathbb Q}
\newcommand{\F}{\mathbb F}
\newcommand{\eps}{\varepsilon}
\begin{document}
% UnsolvedMath ID 30004241; source catalogue code OWR-17137-002
\setcounter{section}{30004240}
\section{Logarithmic Entanglement Enhancement from Partial Delocalization}\label{30004241}
{\small UnsolvedMath ID 30004241 · Mathematical Physics}
\subsection{Definitions and mathematical statement}

\paragraph{Shared notation.}$\mathbb N=\{1,2,\ldots\}$; the standard number fields have their usual meanings. All additional conventions are stated locally.


\paragraph{Definitions and conventions.}
For a bipartite quantum state, entanglement entropy is the von Neumann entropy of a reduced state. In a free-fermion ground state it is computed from the restriction of the one-particle correlation operator to a spatial region. An area law grows like boundary area; logarithmic enhancement adds a factor comparable to the logarithm of region size. Delocalization concerns spatial spread or transport of one-particle states and is not identical to absolute continuity of the spectral measure.

\paragraph{Target retained from the canonical source catalogue.}
\paragraph{Catalogue wording retained for traceability.}
For entanglement entropy in disordered free-fermion systems, how much delocalization of the one-particle eigenvalues and eigenfunctions is needed to produce a logarithmic enhancement of the area law, even when the spectrum is not absolutely continuous?

\paragraph{Source scope and qualification.}
The primary report (PDF pages 21--22) already gives a random-dimer example with logarithmic enhancement without absolutely continuous spectrum. The remaining inherited question asks for how much delocalization and a general criterion; mere existence in one such model is known context, not a freshly verified open subproblem.
\subsection{Short English statement}
How much partial delocalization is enough to create logarithmically enhanced free-fermion entanglement, even without absolutely continuous spectrum?
\subsection{Sources}
\begin{itemize}
\item[\hypertarget{source:30004241:1}{S1}] ULAM.AI and the UnsolvedMath Contributors, UnsolvedMath, record 30004241 (OWR-17137-002). Catalogue curation under CC BY 4.0; definitions and English exposition added locally.
\url{https://huggingface.co/datasets/ulamai/UnsolvedMath}.
\item[\hypertarget{source:30004241:2}{S2}] Many-Body Quantum Systems (2019), Oberwolfach Reports, DOI 10.4171/owr/2019/41, pp. 2541--2603. Primary report accessed 2026-10-11; relevant PDF pages 21, 22.
\url{https://doi.org/10.4171/owr/2019/41}.
\item[\hypertarget{source:30004241:3}{S3}] Muller, Pastur, and Schulte (catalogue-listed context; not independently consulted in this preparation).
\url{https://doi.org/10.1007/s00220-019-03523-3}.
\end{itemize}
\label{end:30004241}
\end{document}
